\documentclass[11pt,fontset=fandol]{ctexart}
\usepackage[a4paper,margin=1in]{geometry}
\usepackage{amsmath,amssymb,amsfonts}
\usepackage{booktabs}
\usepackage{graphicx}
\usepackage{caption}
\usepackage{enumitem}
\usepackage[table]{xcolor}
\usepackage{array}
\usepackage[hidelinks]{hyperref}
\usepackage{xurl}
\usepackage{microtype}

\title{二维 Corridor 半透膜与近起点双 Target 的双峰研究\\\large 新报告（接续 \texttt{grid2d\_rect\_bimodality}）}
\author{}
\date{}

\setlist{nosep}
\setlength{\textfloatsep}{8pt plus 2pt minus 2pt}
\setlength{\floatsep}{6pt plus 2pt minus 2pt}
\setlength{\intextsep}{6pt plus 2pt minus 2pt}
\setlength{\abovecaptionskip}{4pt}
\setlength{\belowcaptionskip}{0pt}
\setlength{\tabcolsep}{5pt}
\renewcommand{\arraystretch}{1.10}

\begin{document}
\maketitle

\section*{结论先行}
本报告在二维长方形域上新增两条研究线，并且与旧报告的 phase 判据、峰谷图式、分解解释保持一致。
\begin{enumerate}
  \item \textbf{有 corridor、单 target：}将原全反射走廊墙替换为半透膜，分别研究对称膜（$\kappa_\uparrow=\kappa_\downarrow$）与非对称膜（$\kappa_\uparrow\neq\kappa_\downarrow$）。
  \item \textbf{无 corridor、双 target：}将一个 target 放在起点附近，另一个保持远端，研究“近目标快峰 + 远目标慢峰”是否稳定形成双峰。
\end{enumerate}

核心结果：
\begin{itemize}
  \item 对称半透膜扫描中，clear-double 在超低泄漏区形成连续带（见图~\ref{fig:sym-phase} 与表~\ref{tab:sym-w16}）；随后进入 weak-double 过渡带，再增大渗透率后转为单峰。
  \item 非对称膜（如 $\kappa_\uparrow=0.002,\kappa_\downarrow=0$）可保持 clear-double，并在窗口分类中显著改变 valley/第二峰的组成权重。
  \item 无 corridor 双 target（近目标放到起点附近）可稳定得到双峰：代表案例 $(b_x,d,\Delta y)=(0.07,2,3)$ 的 near/far 质量约为 $0.326/0.674$、峰位约 $(11,739)$；专用 clear-instance（showcase 最优）为 $(0.12,4,1)$，near/far 质量约 $0.350/0.650$，且 valley/max 约 $0.007$。
\end{itemize}

\section{模型与参数}
\subsection{统一基线}
沿用已有二维长方形框架：
\[
L_x=60,\quad q=0.8,\quad x_s=7,\quad x_t=58,
\]
一阶邻域随机游走 + 惰性停留；边界为反射边界。
单 target corridor 配置使用本轮优化后的锚点：$h=2$、$m=5$、$\delta_{\rm core}=1.00$、$\delta_{\rm open}=0.55$、$b_x=-0.08$。

\subsection{半透膜定义}
原走廊墙是全反射（穿越概率 0）。本报告改为半透膜：
\[
P(\text{穿越膜边})=\kappa,\qquad P(\text{被反射})=1-\kappa.
\]
对称膜：$\kappa_\uparrow=\kappa_\downarrow=\kappa$；
非对称膜：$\kappa_\uparrow\neq\kappa_\downarrow$。

\subsection{分类框架}
保持“峰-谷-峰”窗口分析，并加入新的 4 类标签：
\begin{itemize}
  \item $L$: 是否发生膜泄漏（是否穿过膜边）；
  \item $R$: 是否发生 committor 承诺后回撤（recross）。
\end{itemize}
故有四类：\texttt{L0R0}, \texttt{L0R1}, \texttt{L1R0}, \texttt{L1R1}。

\section{结果 A：corridor + 单 target + 半透膜}
\subsection{对称膜扫描}
\begin{figure}[!htbp]
  \centering
  \includegraphics[width=0.86\linewidth]{figures/membrane_symmetric_phase_map.pdf}
  \caption{对称半透膜相图（纵轴 $W_y$，横轴 $\kappa$，颜色数字为 phase）。}
  \label{fig:sym-phase}
\end{figure}

\begin{figure}[!htbp]
  \centering
  \includegraphics[width=0.86\linewidth]{figures/membrane_symmetric_sep_map.pdf}
  \caption{对称半透膜分离度热图（sep-score）。}
  \label{fig:sym-sep}
\end{figure}

\begin{table}[!htbp]
  \centering
  \footnotesize
  \begin{tabular}{@{}cccccc@{}}
\toprule
$\kappa$ & phase & $t_{p1}$ & $t_{p2}$ & valley/max & sep-score \\
\midrule
0 & 2 & 395 & 1711 & 0.058 & 1.87 \\
0.0003 & 2 & 395 & 1710 & 0.061 & 1.84 \\
0.0005 & 2 & 395 & 1708 & 0.063 & 1.80 \\
0.0008 & 2 & 395 & 1706 & 0.066 & 1.76 \\
0.001 & 2 & 394 & 1704 & 0.069 & 1.72 \\
0.0013 & 2 & 394 & 1702 & 0.072 & 1.68 \\
0.0015 & 2 & 394 & 1700 & 0.074 & 1.64 \\
0.0018 & 2 & 394 & 1698 & 0.077 & 1.59 \\
0.002 & 2 & 394 & 1696 & 0.080 & 1.53 \\
0.0022 & 2 & 393 & 1694 & 0.082 & 1.46 \\
0.0025 & 2 & 393 & 1692 & 0.085 & 1.40 \\
0.0027 & 2 & 393 & 1690 & 0.087 & 1.34 \\
0.003 & 2 & 393 & 1689 & 0.090 & 1.29 \\
0.0032 & 2 & 393 & 1687 & 0.093 & 1.25 \\
0.0035 & 2 & 393 & 1685 & 0.095 & 1.21 \\
0.0037 & 1 & 392 & 1683 & 0.098 & 0.84 \\
0.004 & 1 & 392 & 1681 & 0.100 & 0.84 \\
0.0045 & 1 & 392 & 1677 & 0.105 & 0.82 \\
0.005 & 1 & 392 & 1673 & 0.110 & 0.81 \\
0.006 & 1 & 391 & 1664 & 0.120 & 0.79 \\
0.0075 & 1 & 390 & 1652 & 0.134 & 0.77 \\
0.01 & 1 & 388 & 1629 & 0.157 & 0.73 \\
0.015 & 1 & 386 & 1576 & 0.198 & 0.67 \\
0.02 & 1 & 383 & 1505 & 0.235 & 0.61 \\
0.03 & 0 & - & - & - & 0.00 \\
0.05 & 0 & - & - & - & 0.00 \\
0.1 & 0 & - & - & - & 0.00 \\
0.2 & 0 & - & - & - & 0.00 \\
\bottomrule
\end{tabular}

  \caption{$W_y=16$ 截面的对称膜指标。}
  \label{tab:sym-w16}
\end{table}

在本次参数下（以 $W_y=16$ 截面为例），$\kappa\in[0,0.0035]$ 保持 clear-double；约 $\kappa\in[0.0037,0.02]$ 处于 weak-double；$\kappa\ge 0.03$ 转为单峰，说明“膜泄漏”会持续削弱原走廊机制的时间尺度分离。

\subsection{非对称膜扫描}
\begin{figure}[!htbp]
  \centering
  \includegraphics[width=0.86\linewidth]{figures/membrane_asymmetric_phase_map.pdf}
  \caption{非对称半透膜相图（$\kappa_\uparrow$ vs $\kappa_\downarrow$）。}
  \label{fig:asym-phase}
\end{figure}

\begin{figure}[!htbp]
  \centering
  \includegraphics[width=0.86\linewidth]{figures/membrane_asymmetric_sep_map.pdf}
  \caption{非对称半透膜分离度热图。}
  \label{fig:asym-sep}
\end{figure}

\begin{table}[!htbp]
  \centering
  \footnotesize
  \begin{tabular}{@{}cccccc@{}}
\toprule
$\kappa_{\uparrow}$ & $\kappa_{\downarrow}$ & phase & $t_{p1}$ & $t_{p2}$ & sep-score \\
\midrule
0 & 0.0003 & 2 & 395 & 1711 & 1.85 \\
0.0003 & 0 & 2 & 395 & 1710 & 1.85 \\
0.0005 & 0 & 2 & 395 & 1710 & 1.84 \\
0 & 0.0005 & 2 & 395 & 1710 & 1.84 \\
0.0008 & 0 & 2 & 395 & 1709 & 1.82 \\
0.0005 & 0.0003 & 2 & 395 & 1709 & 1.82 \\
0.0003 & 0.0005 & 2 & 395 & 1709 & 1.82 \\
0 & 0.0008 & 2 & 395 & 1709 & 1.82 \\
0.001 & 0 & 2 & 395 & 1708 & 1.80 \\
0.002 & 0 & 2 & 394 & 1704 & 1.73 \\
\bottomrule
\end{tabular}

  \caption{非对称膜 top cases（按 phase 与 sep 排序）。}
  \label{tab:asym-top}
\end{table}

代表性对比（对称 $\kappa=0.002$ vs 非对称 $\kappa_\uparrow=0.002,\kappa_\downarrow=0$）如下：
\begin{figure}[!htbp]
  \centering
  \includegraphics[width=0.88\linewidth]{figures/membrane_rep_sym_geometry.pdf}
  \caption{对称半透膜代表案例配置图（包含 corridor、膜位置、起点/终点与局部偏置方向）。}
  \label{fig:rep-sym-geo}
\end{figure}

\begin{figure}[!htbp]
  \centering
  \includegraphics[width=0.86\linewidth]{figures/membrane_rep_sym_fpt.pdf}
  \caption{对称半透膜代表案例：总 FPT 与 corridor/outer 通量分解。}
  \label{fig:rep-sym-fpt}
\end{figure}

\begin{figure}[!htbp]
  \centering
  \includegraphics[width=0.88\linewidth]{figures/membrane_rep_asym_geometry.pdf}
  \caption{非对称半透膜代表案例配置图（上/下膜通透率不同）。}
  \label{fig:rep-asym-geo}
\end{figure}

\begin{figure}[!htbp]
  \centering
  \includegraphics[width=0.86\linewidth]{figures/membrane_rep_asym_fpt.pdf}
  \caption{非对称半透膜代表案例：总 FPT 与 corridor/outer 通量分解。}
  \label{fig:rep-asym-fpt}
\end{figure}

\begin{figure}[!htbp]
  \centering
  \includegraphics[width=0.94\linewidth]{figures/membrane_class_window_bars.pdf}
  \caption{峰1/谷/峰2 三窗口中的四类占比（L×R）。}
  \label{fig:membrane-class-bars}
\end{figure}

\begin{figure}[!htbp]
  \centering
  \includegraphics[width=0.90\linewidth]{figures/membrane_class_legend.pdf}
  \caption{L×R 四类代表图例：左图为 2×2 分类定义，右图为代表性窗口主导类别转移示意。}
  \label{fig:membrane-class-legend}
\end{figure}

\begin{table}[!htbp]
  \centering
  \footnotesize
  \begin{tabular}{@{}lccccc@{}}
\toprule
Window & Case & L0R0 & L0R1 & L1R0 & L1R1 \\
\midrule
peak1 & sym & 0.996 & 0.000 & 0.004 & 0.000 \\
peak1 & dir & 0.998 & 0.000 & 0.002 & 0.000 \\
valley & sym & 0.581 & 0.116 & 0.299 & 0.004 \\
valley & dir & 0.655 & 0.121 & 0.224 & 0.001 \\
peak2 & sym & 0.142 & 0.742 & 0.080 & 0.036 \\
peak2 & dir & 0.166 & 0.749 & 0.067 & 0.018 \\
\bottomrule
\end{tabular}

  \caption{四类窗口占比数值表。}
  \label{tab:membrane-window-classes}
\end{table}

\begin{table}[!htbp]
  \centering
  \footnotesize
  \begin{tabular}{@{}lllc c@{}}
\toprule
Case & $q^*$ & Window & dominant class & $w_{\mathrm{L1R1}}$ \\
\midrule
sym & 0.4 & peak1 & L0R0 & 0.000 \\
sym & 0.4 & valley & L0R0 & 0.197 \\
sym & 0.4 & peak2 & L0R1 & 0.104 \\
dir & 0.4 & peak1 & L0R0 & 0.000 \\
dir & 0.4 & valley & L0R0 & 0.196 \\
dir & 0.4 & peak2 & L0R1 & 0.082 \\
sym & 0.5 & peak1 & L0R0 & 0.000 \\
sym & 0.5 & valley & L0R0 & 0.004 \\
sym & 0.5 & peak2 & L0R1 & 0.036 \\
dir & 0.5 & peak1 & L0R0 & 0.000 \\
dir & 0.5 & valley & L0R0 & 0.001 \\
dir & 0.5 & peak2 & L0R1 & 0.018 \\
sym & 0.6 & peak1 & L0R0 & 0.000 \\
sym & 0.6 & valley & L0R0 & 0.002 \\
sym & 0.6 & peak2 & L0R1 & 0.022 \\
dir & 0.6 & peak1 & L0R0 & 0.000 \\
dir & 0.6 & valley & L0R0 & 0.000 \\
dir & 0.6 & peak2 & L0R0 & 0.008 \\
\bottomrule
\end{tabular}

  \caption{$q^*$ 鲁棒性检查（精确分解）：不同承诺阈值下各窗口主导类别与 $w_{\mathrm{L1R1}}$。}
  \label{tab:membrane-qstar-sens}
\end{table}

基于“严格 recross（先达 $\Sigma_{q^*}$ 后回到 $A$）”的精确分解结果可见：peak1 由 \texttt{L0R0} 主导；valley 仍以 \texttt{L0R0} 为主但 \texttt{L1R0} 显著上升；peak2 由 \texttt{L0R1} 主导。说明慢峰主要来自“承诺后回撤”而不一定要求膜泄漏。

\section{结果 B：无 corridor + 双 target（近起点目标）}
本节移除 corridor，仅保留反射边界与全局 $b_x$，并把一个 target 放到起点附近。扫描流程先在候选网格上优化 clear 覆盖率，再从全部候选结果中自动选取 clear-instance（用于“最清晰双峰示例图”）。

\begin{figure}[!htbp]
  \centering
  \includegraphics[width=0.88\linewidth]{figures/two_target_nearstart_phase_map.pdf}
  \caption{无 corridor 双 target：phase map（横轴近目标距离 $d$，纵轴 $b_x$）。}
  \label{fig:tt-near-phase}
\end{figure}

\begin{figure}[!htbp]
  \centering
  \includegraphics[width=0.88\linewidth]{figures/two_target_nearstart_sep_map.pdf}
  \caption{无 corridor 双 target：分离度热图。}
  \label{fig:tt-near-sep}
\end{figure}

\begin{table}[!htbp]
  \centering
  \footnotesize
  \begin{tabular}{@{}cccccc@{}}
\toprule
$b_x$ & clear $d$ interval & best $d$ & phase(best) & $P(\mathrm{near})$ & $P(\mathrm{far})$ \\
\midrule
+0.00 & - & 3 & 0 & 0.892 & 0.108 \\
+0.02 & - & 2 & 0 & 0.603 & 0.397 \\
+0.04 & 3-9 & 3 & 2 & 0.444 & 0.556 \\
+0.05 & 2-10 & 2 & 2 & 0.393 & 0.607 \\
+0.06 & 2-10 & 2 & 2 & 0.355 & 0.645 \\
+0.07 & 2-10 & 2 & 2 & 0.326 & 0.674 \\
+0.08 & 2-10 & 2 & 2 & 0.301 & 0.699 \\
+0.09 & 2-10 & 2 & 2 & 0.281 & 0.719 \\
+0.10 & 2-10 & 2 & 2 & 0.264 & 0.736 \\
+0.11 & 2-10 & 2 & 2 & 0.248 & 0.752 \\
+0.12 & 2-10 & 2 & 2 & 0.235 & 0.765 \\
+0.14 & 2-10 & 2 & 2 & 0.212 & 0.788 \\
+0.16 & 2-10 & 3 & 2 & 0.206 & 0.794 \\
\bottomrule
\end{tabular}

  \caption{按 $b_x$ 汇总的 clear 区间与最优近目标距离。}
  \label{tab:tt-near-summary}
\end{table}

代表案例（$W_y=16,b_x=0.07,d=2$）:
\begin{figure}[!htbp]
  \centering
  \includegraphics[width=0.88\linewidth]{figures/two_target_rep_geometry.pdf}
  \caption{无 corridor 双 target 代表案例配置图（起点、近目标、远目标与全局漂移方向示意）。}
  \label{fig:tt-near-rep-geo}
\end{figure}

\begin{figure}[!htbp]
  \centering
  \includegraphics[width=0.86\linewidth]{figures/two_target_rep_fpt.pdf}
  \caption{无 corridor 双 target 代表案例：总 FPT 及 near/far splitting 分量。}
  \label{fig:tt-near-rep-fpt}
\end{figure}

\begin{figure}[!htbp]
  \centering
  \includegraphics[width=0.94\linewidth]{figures/two_target_rep_doublepeak_example.pdf}
  \caption{无 corridor clear-double 实例图：左图为线性坐标下的尾部增强视图（为突出 valley/peak2，对首峰做上截并以 \texttt{p1*} 标注真值），右图为对数坐标（$y\ge 10^{-6}$）下的双峰结构。}
  \label{fig:tt-near-doublepeak-example}
\end{figure}

\begin{table}[!htbp]
  \centering
  \footnotesize
  \resizebox{\linewidth}{!}{\begin{tabular}{@{}cccccccccccccccc@{}}
\toprule
spec & $b_x$ & $d$ & $\Delta y$ & phase & sep-score & valley/max & $p_{\min}$ & peak-ratio & peak-margin & min-margin & gate-pass & gate-slack & clear-score & showcase & $\varepsilon_{\mathrm{split}}$ \\
\midrule
adaptive\_refine & +0.12 & 4 & 1 & 2 & 2.38 & 0.007 & 0.350 & 0.170 & 0.070 & 0.070 & yes & 0.070 & 1.54 & 1.92 & $1.18\times 10^{-9}$ \\
\bottomrule
\end{tabular}
}
  \caption{无 corridor clear-instance 质量指标（从全部候选扫描中自动选点）。表中同时给出 \texttt{peak-margin}、\texttt{min-margin} 与首道强 gate 审计（\texttt{gate-pass}/\texttt{gate-slack}），避免“刚过阈值但视觉不稳”的伪 clear 点；并附 \texttt{showcase} 用于图表可视化排序一致性。}
  \label{tab:tt-near-clear-instance}
\end{table}

\begin{figure}[!htbp]
  \centering
  \includegraphics[width=0.94\linewidth]{figures/two_target_no_corridor_clear_instance.pdf}
  \caption{无 corridor clear-instance 专用图（全候选扫描中 showcase-score 最大点）：线性尾部增强 + 对数视图并列，确保 valley 与第二峰可视化清晰。}
  \label{fig:tt-near-clear-instance}
\end{figure}

\begin{figure}[!htbp]
  \centering
  \includegraphics[width=0.82\linewidth]{figures/two_target_rep_window_split.pdf}
  \caption{代表案例在 peak1/valley/peak2 窗口的 near/far 质量占比。}
  \label{fig:tt-near-window}
\end{figure}

本轮 clear-instance 的 near/far 总质量、sep-score、valley/max、peak/min 裕量、showcase 分数及 splitting committor 闭环误差 $\varepsilon_{\rm split}$ 见表~\ref{tab:tt-near-clear-instance}。本次在最终排序前加入“稳健门槛筛选”（优先要求 \texttt{peak-margin/min-margin} 正裕量并优先保留深谷约束），确保选出的 no-corridor 实例不仅满足 phase=2 判据，也能在阈值轻微扰动下保持可视化稳定。当前自动选出的 clear-instance 来自细化轮次（spec=\texttt{adaptive\_refine}），参数为 $(b_x,d,\Delta y)=(0.12,4,1)$。其 \texttt{peak-margin}=0.070、\texttt{min-margin}=0.070、\texttt{showcase}=1.92，并通过首道 gate：\texttt{peak>=0.07}、\texttt{min>=0.07}、\texttt{valley<=0.10}；对应审计量 \texttt{gate-pass=yes}、\texttt{gate-slack}=0.070（三条阈值净裕量的最小值），首道 gate 候选数为 168。

\section{方法参考（committor + 2×2 分类）}
\subsection*{phase/sep-score 的可复现定义}
为保证相图可复现，本文直接采用代码中的判据（与 \path{src/vkcore/grid2d/rect_bimodality/cli.py} 一致）：
\paragraph{(1) 双峰探测（\texttt{find\_two\_peaks}）}
令 $\tilde f=\mathrm{smooth}_7(f)$（窗口 7 平滑）。候选峰点需满足局部极大且高于相对阈值
\[
\tilde f(t)>\tilde f(t-1),\quad \tilde f(t)\ge \tilde f(t+1),\quad \tilde f(t)\ge 0.06\,\max \tilde f.
\]
从候选集合中选两峰 $(t_{p1},t_{p2})$ 时，额外要求（代码默认值）：
\[
t_{p2}-t_{p1}\ge 20,\ \mathrm{peak\_balance}\ge 0.08,\ \mathrm{valley\_drop}\ge 0.08,\ \mathrm{minor\_drop}\ge 0.005,
\]
并满足左右主导性容差（\texttt{dominance\_tol}=0.02）。定义
\[
\mathrm{has\_double}:=\mathbf 1\{t_{p1},t_{p2}\ \text{均存在}\}.
\]

\paragraph{(2) one-target 的 phase 定义}
先定义
\[
\mathrm{sep}_{\rm peaks}:=\frac{|t_{p2}-t_{p1}|}{w_{1/2}^{(1)}+w_{1/2}^{(2)}},\quad
\mathrm{valley}/\max:=\frac{h_v}{\max(h_{p1},h_{p2})},\quad
\mathrm{peak\_balance}:=\frac{\min(h_{p1},h_{p2})}{\max(h_{p1},h_{p2})}.
\]
其中 $w_{1/2}^{(k)}$ 为第 $k$ 峰半高宽（由原始 $f$ 计算），$h_{p1},h_{p2},h_v$ 由平滑序列 $\tilde f$ 在峰/谷位置读取。
分类为
\[
\mathrm{phase}_{\rm OT}=
\begin{cases}
2,& \mathrm{has\_double}=1,\ \mathrm{sep}_{\rm peaks}\ge 1.0,\ \mathrm{valley}/\max\le 0.35,\ \mathrm{peak\_balance}\ge 0.15,\\
1,& \mathrm{has\_double}=1\ \text{但不满足上式},\\
0,& \mathrm{has\_double}=0.
\end{cases}
\]

\paragraph{(3) two-target 的 phase 定义}
先定义分量模态时间与分离度
\[
t_{\rm mode}^{\rm near}:=\arg\max_t f_{\rm near}(t),\quad
t_{\rm mode}^{\rm far}:=\arg\max_t f_{\rm far}(t),
\]
\[
\mathrm{sep}_{\rm mode}:=\frac{|t_{\rm mode}^{\rm far}-t_{\rm mode}^{\rm near}|}{w_{1/2}^{\rm far}+w_{1/2}^{\rm near}}.
\]
质量占比由分量总质量归一化得到
\[
P_{\rm near}:=\frac{\sum_t f_{\rm near}(t)}{\sum_t f_{\rm any}(t)},\qquad
P_{\rm far}:=\frac{\sum_t f_{\rm far}(t)}{\sum_t f_{\rm any}(t)}.
\]
并在总分布上定义
\[
\mathrm{peak\_ratio}:=\frac{\min(h_{p1},h_{p2})}{\max(h_{p1},h_{p2})},\qquad
\mathrm{valley}/\max:=\frac{h_v}{\max(h_{p1},h_{p2})}.
\]
分类为
\[
\mathrm{phase}_{\rm TT}=
\begin{cases}
2,& \mathrm{has\_double}=1,\ \mathrm{sep}_{\rm mode}\ge 1.0,\ \min(P_{\rm near},P_{\rm far})\ge 0.15,\\
& \mathrm{valley}/\max\le 0.35,\ \mathrm{peak\_ratio}\ge 0.10,\ (\mathrm{peak\_ratio}-\mathrm{valley}/\max)\ge 0.01,\\
1,& \mathrm{has\_double}=1\ \text{但不满足上式},\\
0,& \mathrm{has\_double}=0.
\end{cases}
\]
因此：\textbf{phase=1} 的数学含义是“检测到双峰，但分离度/平衡度/谷深等条件至少有一项不达 clear-double 阈值”。

\subsection*{committor 的数学机理（详细版）}
记 hitting time 为 $\tau_C:=\inf\{t\ge 0:\,X_t\in C\}$。在离散马尔可夫链转移矩阵 $P$ 上，设 $A$ 为起点盆（本报告用起点附近 strip），$B$ 为目标吸收集合，定义
\[
q_i := \mathbb P_i(\tau_B < \tau_A).
\]
它满足离散 Dirichlet 边值问题
\[
q_i=0\ (i\in A),\qquad q_i=1\ (i\in B),\qquad
q_i=\sum_j P_{ij} q_j\ (i\notin A\cup B).
\]
令 $U=\Omega\setminus(A\cup B)$，按 $(U,A,B)$ 重排后可写成线性系统
\[
(I-P_{UU})\,q_U=P_{UB}\,\mathbf 1.
\]
在本问题中（存在吸收且从 $U$ 终将离开）该系统解唯一，因此 $q$ 不是经验坐标，而是由动力学唯一确定的最优进度坐标。我们取 $q^*=0.5$ 定义承诺面
\[
\Sigma_{q^*}:=\{i:\ q_i\ge q^*\}.
\]
此外，在谱半径 $\rho(P_{UU})<1$ 时有 Neumann 展开
\[
q_U=(I-P_{UU})^{-1}P_{UB}\mathbf 1
=\sum_{n\ge 0}P_{UU}^n P_{UB}\mathbf 1.
\]
这意味着 $q_U$ 可以解释为“在内部区反复游走 $n$ 步后吸收到目标”的总和，从线性系统角度直接对应到后文的 renewal 失败-重试图景：早峰偏向小 $n$，慢峰来自大 $n$ 的累计权重。

\subsection*{L×R 事件定义与 recross 计数}
对轨迹 $X_t$（直到 $\tau_B$）定义膜泄漏次数
\[
N_{\rm mem}:=\sum_{t=0}^{\tau_B-1}\mathbf 1\!\{(X_t,X_{t+1})\in\mathcal E_{\rm mem}\},
\]
以及承诺后回撤次数
\[
N_{\rm ret}:=\sum_{t=0}^{\tau_B-1}\mathbf 1\!\{X_t\in\Sigma_{q^*},\,X_{t+1}\in A\}.
\]
二值分类变量取
\[
L=\mathbf 1\{N_{\rm mem}\ge 1\},\qquad R=\mathbf 1\{N_{\rm ret}\ge 1\},
\]
从而得到四类
\[
(\!L,R\!)\in\{(0,0),(0,1),(1,0),(1,1)\},
\]
对应 \texttt{L0R0/L0R1/L1R0/L1R1}（图~\ref{fig:membrane-class-legend}）。

\subsection*{与双峰/谷对应的 renewal 展开}
将首达分布记为 $f(t)$。若把“一次尝试直接成功”记为核 $g(t)$，把“先失败并回到再生集合后继续”记为核 $h(t)$，则
\[
f(t)=g(t)+(h*f)(t)=\sum_{k\ge 0}(h^{*k}*g)(t).
\]
其中 $k$ 可解释为失败尝试次数。双峰出现的条件就是 $k=0$ 与 $k\ge1$ 分量在时间尺度和质量上都可分离。窗口占比
\[
w_{\ell r}(W)=\frac{\sum_{t\in W} f_{\ell r}(t)}{\sum_{t\in W} f(t)}
\]
给出可视化证据：peak1 主要由 \texttt{L0R0} 托起；valley 仍以 \texttt{L0R0} 为主、但 \texttt{L1R0} 明显上升；peak2 则以 \texttt{L0R1} 主导，\texttt{L1R1} 仅占次要比例（图~\ref{fig:membrane-class-bars}，表~\ref{tab:membrane-window-classes}）。

\subsection*{从线性系统到首达核的闭环公式}
在无 corridor 双 target 场景，把非吸收内部态记为 $U$，起点内部基向量记为 $e_s$，并记
\[
Q:=P_{UU},\qquad R_{\rm near}:=P_{U,\partial_{\rm near}},\qquad R_{\rm far}:=P_{U,\partial_{\rm far}}.
\]
则首达 near/far 的 pmf 可写为
\[
f_{\rm near}(t)=e_s^\top Q^{t-1}R_{\rm near}\mathbf 1,\qquad
f_{\rm far}(t)=e_s^\top Q^{t-1}R_{\rm far}\mathbf 1,\qquad t\ge 1.
\]
相应总质量为
\[
P_{\rm far}=\sum_{t\ge 1}f_{\rm far}(t)=e_s^\top(I-Q)^{-1}R_{\rm far}\mathbf 1.
\]
这条公式把“时域首达曲线分解”与“线性系统解”直接连通，是后续 splitting committor 闭环的核心。

\subsection*{条件化核（Doob 变换）与峰谷转移解释}
在内部态且 $q_{\rm far}(i)>0$ 时，定义
\[
\widetilde P^{\rm far}_{ij}:=\frac{P_{ij}\,q_{\rm far}(j)}{q_{\rm far}(i)}.
\]
这就是以 $h=q_{\rm far}$ 对原转移核做的 Doob $h$-transform，对应“条件在最终先到 far 再到 near”下的精确动力学；near 分支则由互补的 $q_{\rm near}=1-q_{\rm far}$ 给出。于是 peak1 到 peak2 的窗口转移可解释为两类条件化轨道测度的再分配：早期质量主要来自短程成功轨道，而晚期质量来自多次离开并回到低-$q_{\rm far}$ 区域后的最终吸收。这一视角把观测到的 renewal 分解与可计算的条件化动力学直接对应起来。

\subsection*{谱分解视角：快慢双时标从何而来}
在内部态子核 $Q$ 上，若写成（广义）谱展开
\[
Q^{t-1}=\sum_{k}\lambda_k^{\,t-1} r_k\ell_k^\top,\qquad |\lambda_1|>|\lambda_2|\ge\cdots,
\]
则
\[
f_{\rm near}(t)=\sum_k c_k^{\rm near}\lambda_k^{\,t-1},\qquad
f_{\rm far}(t)=\sum_k c_k^{\rm far}\lambda_k^{\,t-1}.
\]
较小模值模态在早期快速衰减，主模 $\lambda_1$ 控制慢尾和晚峰。当 near/far 两分量在模态投影 $c_k^{\rm near},c_k^{\rm far}$ 上权重不同，就会出现“早期 near 主导、晚期 far 主导”的窗口迁移。这与 renewal 分解一致：peak1 对应短停留（小循环）贡献，peak2 对应准稳态主模累积释放。等价地看条件风险率
\[
h_{\rm far}(t):=\frac{f_{\rm far}(t)}{S(t-1)},
\]
其随时间上升就意味着远端吸收通道在晚窗口逐步占优。

\subsection*{双 target splitting committor 闭环}
对同一系统，定义
\[
q_{\rm far}(i):=\mathbb P_i(\tau_{\rm far}<\tau_{\rm near}),\qquad q_{\rm near}(i)=1-q_{\rm far}(i).
\]
其边值问题为
\[
q_{\rm far}=0\ \text{on near},\qquad q_{\rm far}=1\ \text{on far},\qquad
q_{\rm far}(i)=\sum_j P_{ij}\,q_{\rm far}(j)\ \text{on interior}.
\]
因此一级分解 $f=f_{\rm near}+f_{\rm far}$ 与 committor 可直接闭环：$q_{\rm far}(x_s)$ 应与 $P_{\rm far}$ 一致。定义
\[
\varepsilon_{\rm split}:=\left|q_{\rm far}(x_s)-\sum_{t\ge 1}f_{\rm far}(t)\right|
=\left|q_{\rm far}(x_s)-P_{\rm far}\right|.
\]
进一步地，令 $\tau:=\tau_{\rm near}\wedge\tau_{\rm far}$，并定义
\[
M_t:=q_{\rm far}(X_{t\wedge\tau}).
\]
由于 $q_{\rm far}$ 在 interior 上满足调和关系、在 near/far 边界上取 $0/1$，$M_t$ 是有界鞅。由可选停时定理可得
\[
q_{\rm far}(x_s)=\mathbb E[M_\tau]
=\mathbb P_{x_s}(\tau_{\rm far}<\tau_{\rm near})
=P_{\rm far}.
\]
因此 $\varepsilon_{\rm split}$ 理论上应为 $0$；数值实现中它的量级应与 Bellman 残差处于同一误差层级，可直接作为机制闭环与数值稳定性的联合质检量。
\begin{table}[!htbp]
  \centering
  \footnotesize
  \resizebox{\linewidth}{!}{\begin{tabular}{@{}lcccccc@{}}
\toprule
case & $q_{\mathrm{far}}(x_s)$ & $P_{\mathrm{far}}$ & $\varepsilon_{\mathrm{split}}$ & $\|r_{\mathrm{int}}\|_\infty$ & $\|r_{\mathrm{near}}\|_\infty$ & $\|r_{\mathrm{far}}\|_\infty$ \\
\midrule
representative & 0.676009 & 0.676009 & $1.14\times 10^{-9}$ & $9.93\times 10^{-12}$ & $0$ & $0$ \\
clear-instance & 0.650264 & 0.650264 & $1.18\times 10^{-9}$ & $9.93\times 10^{-12}$ & $0$ & $0$ \\
\bottomrule
\end{tabular}
}
  \caption{no-corridor 双 target 的 splitting committor 闭环质检（代表案例 + 专用 clear-instance）：同时报告 $q_{\rm far}(x_s)$、$P_{\rm far}$、$\varepsilon_{\rm split}$ 与内点/边界残差。}
  \label{tab:tt-splitting-qc}
\end{table}

表~\ref{tab:tt-splitting-qc} 显示，代表路径与 clear-instance 路径都满足 $q_{\rm far}(x_s)\approx P_{\rm far}$，且 $\varepsilon_{\rm split}$ 分别在 $10^{-8}$ 与 $10^{-9}$ 量级，内点 Bellman 残差保持在 $10^{-11}$ 量级。说明“快/慢峰由 near/far 吸收竞争产生”的解释在两条路径上都与精确分解一致。后续迭代中把 $\varepsilon_{\rm split}$ 与 committor 在内点的 Bellman 残差、near/far 边界残差一起作为自动质检量，可持续监控图表解释与数值分解是否漂移。

\subsection*{本报告中的 committor 实现细节}
\begin{enumerate}
  \item 在代表案例上先构建精确转移矩阵（已包含懒惰项、偏置、膜穿透率和反射边界）。
  \item 取 $A=\{x\le x_s+1,\ |y-y_{\rm mid}|\le 1\}$，$B=\{\text{target}\}$，并用固定点迭代求解 $q$（迭代终止阈值 $10^{-11}$）；同时输出内点 Bellman 残差与边界残差作为数值质检。
  \item 对 L×R 采用扩展状态空间 $(x,\ell,s,r)$ 的\textbf{精确递推}而非 MC：其中 $\ell$ 记录是否发生膜泄漏，$s$ 记录是否访问过 $\Sigma_{q^*}$，$r$ 记录是否在访问 $\Sigma_{q^*}$ 后返回 $A$。
  \item 默认报告 $q^*=0.5$，并额外输出 $q^*\in\{0.4,0.5,0.6\}$ 的鲁棒性表（表~\ref{tab:membrane-qstar-sens}）。
  \item 在双 target 无 corridor 场景中，一级分解始终采用严格恒等式 $f=f_{\rm near}+f_{\rm far}$，保证“双峰来源”可直接追溯。
\end{enumerate}

\section{总结与后续优化方向（committor 导向）}
\subsection*{实现总结（本轮优化后）}
\begin{enumerate}
  \item corridor 半透膜：对称/非对称代表案例均保持 clear-double；在严格 recross 定义下，peak1 主要是 \texttt{L0R0}，而慢峰由 \texttt{L0R1} 抬升，valley 中 \texttt{L1R0} 比例明显增加。
  \item 无 corridor 双 target：新增“全候选扫描 $\rightarrow$ 自适应局部细化 $\rightarrow$ 稳健门槛筛选 $\rightarrow$ clear-instance 排序”流程，并落地专用图/表（图~\ref{fig:tt-near-clear-instance}，表~\ref{tab:tt-near-clear-instance}）；当前 clear-instance 为 $(b_x,d,\Delta y)=(0.12,4,1)$，在 five-constraint 下保持正裕量，并通过首道 gate（\texttt{peak>=0.07}、\texttt{min>=0.07}、\texttt{valley<=0.10}）；新增强 gate 审计字段（\texttt{gate-pass/gate-slack}）和候选数统计（首道 gate 候选 168 个），同时保留 clear-instance 图完整性硬校验（文件存在且大小 $>2$KB），并把图路径/大小写入 `summary.json` 供持续巡检。
  \item 方法层面：committor 从“术语”落到了可计算、可复核流程；splitting committor 在代表点与 clear-instance 上都与吸收质量形成数值闭环，并把 $\varepsilon_{\rm split}$ 与 committor 残差共同纳入自动质检（当前代表点 $\varepsilon_{\rm split}$ 量级 $\sim10^{-8}$，clear-instance 量级 $\sim10^{-9}$）。
  \item 工程执行层面：本轮本地完整回归（codegen + 中英文构建）单轮耗时约 $70\,$s（约 1.2 分钟），满足“本地 $<10$ 分钟优先”条件；后续仅在扫描规模扩大到预计单轮 $>20$ 分钟时切换 Isambard。
\end{enumerate}

\subsection*{后续方向（全部围绕 committor 落地）}
\begin{enumerate}
  \item 将当前“代表案例精确分解”扩展为“整张相图上的精确分解”，并同步输出 clear-instance 的阈值裕量热图（而不只给单点表格）。
  \item 做 $q^*$ 鲁棒性扫描（如 $0.4/0.5/0.6$）并报告峰谷分类稳定性，避免阈值依赖质疑。
  \item 在双 target 场景系统化输出 splitting committor 场与窗口统计
  \[
  q_{\rm far}(x)=\mathbb P_x(\tau_{\rm far}<\tau_{\rm near}),
  \]
  并把 $\varepsilon_{\rm split}=|q_{\rm far}(x_s)-P_{\rm far}|$（同时覆盖代表点与 clear-instance 选点）与两项阈值裕量指标（\texttt{min-margin} 与 \texttt{gate-slack}）共同作为自动质检项。
  \item 在膜模型中比较“上下非对称”与“方向非对称（内外穿透率不同）”两类破缺，并统一投影到 committor 几何上对比。
  \item 在 keepalive 循环中把“codegen + 中英文构建可通过”作为一等优化目标，并固定执行分流策略：本地小轮次（$<10$ 分钟）优先直跑；大规模扫描（预计单轮 $>20$ 分钟）走 isbard doctor/auth + submit/status/fetch；远端失败时先记录原因并回退最小本地 smoke。
  \item 所有远端轮次统一输出运行元信息（\texttt{JOB\_ID}、\texttt{REMOTE\_DIR}、关键产物路径与下一步动作），保证总结可追溯。
\end{enumerate}

\section*{复现命令}
\begin{verbatim}
cd reports/grid2d_membrane_near_target
../../.venv/bin/python code/membrane_near_target_report.py
latexmk -xelatex -interaction=nonstopmode -halt-on-error \
  -auxdir=build -emulate-aux-dir grid2d_membrane_near_target_cn.tex
\end{verbatim}

\end{document}
